\section*{NeurIPS Paper Checklist}

\begin{enumerate}

\item {\bf Claims}
    \item[] Question: Do the main claims made in the abstract and introduction accurately reflect the paper's contributions and scope?
    \item[] Answer: \answerYes{}
    \item[] Justification: The abstract and introduction state four concrete contributions: (1) a formal decomposition of performance gaps into sample-size and subset components, (2) the SOKED algorithm, (3) the ordering-reversal failure mode, and (4) empirical demonstrations across nine datasets. All four are delivered in the paper.

\item {\bf Limitations}
    \item[] Question: Does the paper discuss the limitations of the work performed by the authors?
    \item[] Answer: \answerYes{}
    \item[] Justification: The Conclusion section notes that SOKED matches at the minimum training size (the most conservative choice), and that future work is needed for power analysis to guide the number of folds and downsampling seeds. The decomposition is operational rather than causal, and this is acknowledged in Section~3.1.

\item {\bf Theory assumptions and proofs}
    \item[] Question: For each theoretical result, does the paper provide the full set of assumptions and a complete (and correct) proof?
    \item[] Answer: \answerYes{}
    \item[] Justification: Lemma~1 and Proposition~1 state all assumptions explicitly. Full proofs are provided in Appendix~A.

\item {\bf Experimental result reproducibility}
    \item[] Question: Does the paper fully disclose all the information needed to reproduce the main experimental results of the paper to the extent that it affects the main claims and/or conclusions of the paper (regardless of whether the code and data are provided or not)?
    \item[] Answer: \answerYes{}
    \item[] Justification: Section~3.3 specifies the learner, number of CV folds, number of downsampling seeds, and matched training size $n_0$ for each dataset. SOKED is implemented in the publicly available \texttt{mlr3resampling} R package.

\item {\bf Open access to data and code}
    \item[] Question: Does the paper provide open access to the data and code, with sufficient instructions to faithfully reproduce the main experimental results, as described in supplemental material?
    \item[] Answer: \answerYes{}
    \item[] Justification: The SOKED framework is implemented in the \texttt{mlr3resampling} R package (available on CRAN). STL-10, UCI Waveform, and NSCH are publicly available datasets. The workers' compensation dataset is publicly available via the cited source. Experiment scripts are available in the supplemental material.

\item {\bf Experimental setting/details}
    \item[] Question: Does the paper specify all the training and test details (e.g., data splits, hyperparameters, how they were chosen, type of optimizer) necessary to understand the results?
    \item[] Answer: \answerYes{}
    \item[] Justification: Section~3.3 reports the learner, number of CV folds ($B=5$ or $B=10$), number of downsampling seeds ($R=1$ or $R=100$), and matched training size for each experiment. Learners use default hyperparameters from \texttt{mlr3}.

\item {\bf Experiment statistical significance}
    \item[] Question: Does the paper report error bars suitably and correctly defined or other appropriate information about the statistical significance of the experiments?
    \item[] Answer: \answerYes{}
    \item[] Justification: All figures report mean $\pm$ SD or mean $\pm$ SE across CV folds. Paired $t$-test $p$-values comparing Same vs.\ Other at matched training size are reported for every dataset, with the test procedure described in Section~3.2.

\item {\bf Experiments compute resources}
    \item[] Question: For each experiment, does the paper provide sufficient information on the computer resources (type of compute workers, memory, time of execution) needed to reproduce the experiments?
    \item[] Answer: \answerYes{}
    \item[] Justification: Figure~2 (Panel B) reports per-run training time for each learner on the traffic dataset. Experiments were run on a SLURM HPC cluster; individual job resource requirements (walltime 2h, memory 24GB) are specified in the supplemental scripts.

\item {\bf Code of ethics}
    \item[] Question: Does the research conducted in the paper conform, in every respect, with the NeurIPS Code of Ethics?
    \item[] Answer: \answerYes{}
    \item[] Justification: The paper uses only publicly available datasets and introduces a diagnostic methodology. No data was collected from human subjects by the authors. The NSCH autism data is an anonymized public survey; the workers' compensation data is synthetic.

\item {\bf Broader impacts}
    \item[] Question: Does the paper discuss both potential positive societal impacts and negative societal impacts of the work performed?
    \item[] Answer: \answerYes{}
    \item[] Justification: The paper explicitly discusses fairness auditing as a use case (Related Work and workers' compensation experiment): SOKED separates sample-size effects from genuine distributional differences, which has positive implications for equitable model evaluation across demographic groups. No significant negative societal impact is anticipated for a model evaluation methodology.

\item {\bf Safeguards}
    \item[] Question: Does the paper describe safeguards that have been put in place for responsible release of data or models that have a high risk for misuse?
    \item[] Answer: \answerNA{}
    \item[] Justification: The paper introduces a cross-validation methodology and releases no new datasets or generative models that pose misuse risks.

\item {\bf Licenses for existing assets}
    \item[] Question: Are the creators or original owners of assets used in the paper properly credited and are the license and terms of use explicitly mentioned and properly respected?
    \item[] Answer: \answerYes{}
    \item[] Justification: All datasets are cited: STL-10 \citep{coates2011stl10}, UCI Waveform \citep{breiman1984classification}, workers' compensation \citep{Ali2020}, NSCH (U.S.\ Census Bureau public-use file), traffic and qSIP (cited co-authors). The \texttt{mlr3resampling} package is MIT-licensed.

\item {\bf New assets}
    \item[] Question: Are new assets introduced in the paper well documented and is the documentation provided alongside the assets?
    \item[] Answer: \answerYes{}
    \item[] Justification: The primary new asset is the SOKED implementation in \texttt{mlr3resampling}, which is documented on CRAN with vignettes and a function reference. The \texttt{ResamplingSameOtherSizesCV} class and \texttt{pvalue\_downsample()} function implement the methods described in Sections~3.1--3.2.

\item {\bf Crowdsourcing and research with human subjects}
    \item[] Question: Does the paper include the full text of instructions given to participants and screenshots, if applicable, as well as details about compensation?
    \item[] Answer: \answerNA{}
    \item[] Justification: The paper does not involve crowdsourcing or direct research with human subjects. The NSCH dataset is a pre-existing public survey conducted by the U.S.\ Census Bureau.

\item {\bf Institutional review board (IRB) approvals or equivalent for research with human subjects}
    \item[] Question: Does the paper describe potential risks incurred by study participants and whether IRB approvals were obtained?
    \item[] Answer: \answerNA{}
    \item[] Justification: No human subjects research was conducted by the authors. The NSCH autism dataset is a de-identified public-use file; the workers' compensation dataset is synthetic.

\item {\bf Declaration of LLM usage}
    \item[] Question: Does the paper describe the usage of LLMs if it is an important, original, or non-standard component of the core methods in this research?
    \item[] Answer: \answerNA{}
    \item[] Justification: LLMs were not used as part of the core methodology, algorithms, or experimental design.

\end{enumerate}
